
%% "Patch" to keep the text widths' of Metropolis similar to that of Madrid.
%\setbeamersize{text margin left=15pt, text margin right=15pt}
%\makeatletter
%\setbeamertemplate{title page}{
%\centering
%  \begin{minipage}[b][\paperheight]{.9\textwidth}
%    \ifx\inserttitlegraphic\@empty\else\usebeamertemplate*{title graphic}\fi
%    \vfill%
%    \ifx\inserttitle\@empty\else\usebeamertemplate*{title}\fi
%    \ifx\insertsubtitle\@empty\else\usebeamertemplate*{subtitle}\fi
%    \usebeamertemplate*{title separator}
%    \ifx\beamer@shortauthor\@empty\else\usebeamertemplate*{author}\fi
%    \ifx\insertdate\@empty\else\usebeamertemplate*{date}\fi
%    \ifx\insertinstitute\@empty\else\usebeamertemplate*{institute}\fi
%    \vfill
%    \vspace*{1mm}
%  \end{minipage}
%}
%\makeatother

% Modify the text widths' of section title pages. See `beamerinnerthememetropolis.dtx`.
% `\paperheight` requires adjustment after modifying the frame content placement. Not sure why.
\makeatletter
\setbeamertemplate{section page}{
  \centering
  \begin{minipage}[c][.95\paperheight]{.9\textwidth}
    \raggedright
    \usebeamercolor[fg]{section title}
    \usebeamerfont{section title}
    \insertsectionhead\\[-1ex]
    \usebeamertemplate*{progress bar in section page}
    \par
    \ifx\insertsubsectionhead\@empty\else%
      \usebeamercolor[fg]{subsection title}%
      \usebeamerfont{subsection title}%
      \insertsubsectionhead
    \fi
  \end{minipage}
  \par
  \vspace{\baselineskip}
}
\makeatother

% Make frametitles larger and place closer to center. See `beamerinnerthememetropolis.dtx`.
\setbeamerfont{frametitle}{size=\Large}
\makeatletter
% Theme default adds paddings of 2.2ex to all the four sides of frame titles
\newlength{\metropolis@frametitle@toppadding}
\newlength{\metropolis@frametitle@bottompadding}
\setlength{\metropolis@frametitle@toppadding}{3.5ex} 
\setlength{\metropolis@frametitle@bottompadding}{0ex} 
\setlength{\metropolis@frametitle@padding}{4ex} % Horizontal padding to left
\renewcommand{\metropolis@frametitlestrut@start}{
  \rule{0pt}{\metropolis@frametitle@toppadding +%
    \totalheightof{%
      \ifcsdef{metropolis@frametitleformat}{\metropolis@frametitleformat X}{X}%
    }%
  }%
}
\newcommand*\getlength[1]{\number#1}
\renewcommand{\metropolis@frametitlestrut@end}{%
  \ifnum\getlength{\metropolis@frametitle@bottompadding}>0%
    \rule[-\metropolis@frametitle@bottompadding]{0pt}{\metropolis@frametitle@bottompadding}%
  \fi%
}
\makeatother

% Move up the frame content a bit. See:
% https://github.com/matze/mtheme/blob/master/source/beamerinnerthememetropolis.dtx#L509-L523
% https://tex.stackexchange.com/questions/247826/beamer-full-vertical-centering
\makeatletter
\define@key{beamerframe}{c}[true]{% centered
  \beamer@frametopskip=0pt plus 0.5fill\relax% Default is `1fill`
  \beamer@framebottomskip=0pt plus 1.5fill\relax%
  \beamer@frametopskipautobreak=0pt plus .4\paperheight\relax%
  \beamer@framebottomskipautobreak=0pt plus .6\paperheight\relax%
  \def\beamer@initfirstlineunskip{}%
}
\makeatother



%% Beamer options.
\setbeamercovered{dynamic}
\setbeamercovered{invisible}
\beamertemplatenavigationsymbolsempty
%\setbeamertemplate{caption}{unnumbered}
\setbeamertemplate{footline}{}
%\setbeameroption{show notes on second screen}

%% Customize styles inside itemize environments
\addtolength{\leftmargini}{-\labelsep}
\setbeamerfont{itemize/enumerate subbody}{size=\normalsize} %to set the body size
%\setbeamertemplate{frametitle}[default][center]
\newcommand{\smallerblacktriangleright}{\raisebox{.2\height}{\scalebox{.75}{$\blacktriangleright$}}}
%\setbeamertemplate{itemize item}{\normalsize\raise1.25pt\hbox{\donotcoloroutermaths$\blacktriangleright$}}  %to set the symbol size
\setbeamertemplate{itemize item}{$\smallerblacktriangleright$}  %to set the symbol size



%% Custom environments
\newcommand{\defineTightSpacing}{%
	\setlength{\abovedisplayskip}{.3\baselineskip}%
	\setlength{\belowdisplayskip}{.3\baselineskip}%
}
\newenvironment{tightEquation}{%
	\defineTightSpacing%
	\begin{equation}
}{
	\end{equation} \ignorespacesafterend
}
\makeatletter
\newenvironment{tightEquation*}{%
	\defineTightSpacing%
	\begin{equation*}
}{
	\end{equation*} \ignorespacesafterend
}

\newcommand{\defineWithinItemizeSpacing}{%
	\setlength{\abovedisplayskip}{.3\baselineskip}%
	\setlength{\belowdisplayskip}{.3\baselineskip}%
}
\newenvironment{itemizedEquation}{%
	\defineWithinItemizeSpacing%
	\begin{equation}
}{
	\end{equation} \ignorespacesafterend
}
\makeatletter
\newenvironment{itemizedEquation*}{%
	\defineWithinItemizeSpacing%
	\begin{equation*}
}{
	\end{equation*} \ignorespacesafterend
}

\newenvironment{indented}[1][3]{%
	\hspace*{#1ex}\begin{minipage}{\dimexpr\textwidth-#1ex} 
	}{
	\end{minipage}
}
\newenvironment{looseItemize}{%
  \begin{itemize}
  \addtolength\itemsep{.3\baselineskip}
}{
  \end{itemize}
}

\newenvironment{tightItemize}[1][]{%
  \vspace{-.3\baselineskip}%
  \begin{itemize}[#1]
  \addtolength\itemsep{-.1\baselineskip}
}{
  \end{itemize}
}
\newenvironment{tightEnumerate}[1][1.]{%
  \vspace{-.3\baselineskip}%
  \begin{enumerate}[#1]
  \addtolength\itemsep{-.1\baselineskip}
}{
  \end{enumerate}
}



% Macros

%% Beamer macros
\newcommand{\titleSmallCap}[1]{{\fontsize{12}{14}\selectfont  \MakeUppercase{#1}}}
\newcommand{\smallCap}[1]{{\small \MakeUppercase{#1}}}

%% Utility macros
\newcommand{\noteBullet}{\hspace*{.75em}\textcolor{themecolor}{$\blacktriangleright$}\ }
\newcommand{\tightSpacing}[2][.55ex]{%
  \begingroup
    \fontdimen2\font=#1
    #2%
  \endgroup
}

%% General math macros
\newcommand{\spacedColon}{\mkern .5mu : \mkern 1mu}
\newcommand{\spacedEq}{\mkern 1mu = \mkern 1.5mu}
\newcommand{\given}{\thinnerspace | \thinnerspace}
\newcommand{\divby}{\thinnerspace /}
\newcommand{\defeq}{\vcentcolon =} % `\coloneqq` and `\vcentcolon` requires the `mathtools` package: https://tex.stackexchange.com/questions/194344/symbol-for-definition
\newcommand{\diff}{\operatorname{\mathrm{d}}\!{}}
\newcommand{\sign}{\mathrm{sign}}
\newcommand{\diag}{\mathrm{diag}}
\newcommand{\medcap}{\mathbin{\mathsmaller{\bigcap}}}
\DeclareMathOperator*{\argmin}{argmin}
\DeclareMathOperator*{\argmax}{argmax}
\newcommand{\elemwiseProd}{\odot}
\renewcommand{\complement}{{\raisebox{1.5pt}{\scriptsize $\mathsf{c}$}}}
\newcommand{\transpose}{\text{\raisebox{.5ex}{$\intercal$}}}
\newcommand{\lowerscript}[1]{\raisebox{-2pt}{\scriptsize $#1$}}
\newcommand{\yesnumber}{\addtocounter{equation}{1}\tag{\theequation}}
\newcommand{\thinnerspace}{\mskip.5\thinmuskip}
\newcommand{\thinnestspace}{\mskip.25\thinmuskip}
\newcommand{\negthinnerspace}{\mskip-.5\thinmuskip}
\newcommand{\spaceBeforePartial}{\mskip\thinmuskip}
\newcommand{\scriptsumi}{{\scriptsize \sum_i}} % Subscript needs to be within bracket, hence the hard coding
\newcommand{\placeholder}{\thinnerspace \cdot \thinnerspace}
\newcommand{\tightTo}{\mskip-.4\thinmuskip \to \mskip-.4\thinmuskip}

%% Probability / statistics macros
\newcommand{\probability}{\mathbb{P}}
\DeclareMathOperator{\expectation}{\mathbb{E}}
\newcommand{\variance}{\mathrm{Var}}
\DeclareMathOperator*{\covariance}{Cov}
\newcommand{\indicator}{\mathds{1}}
\newcommand{\eqDistribution}{\mathrel{\raisebox{-.2ex}{$\overset{\scalebox{.6}{$\, d$}}{=}$}}}
\newcommand{\iidSim}{\mathrel{\raisebox{-.3ex}{$\overset{\text{i.i.d.}}{\sim}$}}}
\newcommand{\unifDist}{\mathrm{Unif}}
\newcommand{\normalDist}{\mathcal{N}}
\newcommand{\betaDist}{\mathrm{Beta}}
\newcommand{\gammaDist}{\mathrm{Gamma}}
\newcommand{\mle}[1]{\widehat{#1}_{\textrm{mle}}}
\newcommand{\map}[1]{\widehat{#1}_{\textrm{map}}}
\newcommand{\empiricalVar}{\widehat{\variance}}
\newcommand{\kldivergence}{D_{\mathrm{KL}}}
\newcommand{\truthSub}{\mathrm{tru}}

%% Variable macros / Aliases
\newcommand{\nPred}{p}
\DeclareMathOperator{\density}{\pi}
\newcommand{\likelihood}{L}
\newcommand{\infoMat}{\mathcal{I}}
\newcommand{\by}{\bm{y}}
\newcommand{\bx}{\bm{x}}
\newcommand{\bX}{\bm{X}}
\newcommand{\bmu}{\bm{\mu}}
\newcommand{\bbeta}{\bm{\beta}}
\newcommand{\btheta}{\bm{\theta}}
\newcommand{\bPhi}{\bm{\Phi}}
\newcommand{\bSigma}{\bm{\Sigma}}
\newcommand{\minus}{\scalebox{1.0}[1.0]{{\scriptsize -}}}
\newcommand{\Id}{\bm{I}}


\newcommand{\mcmc}{{\small MCMC}}
\newcommand{\Mcmc}{M{\small CMC}}
\newcommand{\bda}{\smallCap{bda}}
